% common preamble of the standalone TikZ figures (cropped with the preview package)
\documentclass[11pt]{article}
\usepackage[T1]{fontenc}
\usepackage{lmodern}
\usepackage{amsmath}
\usepackage{xcolor}
\usepackage{tikz}
\usetikzlibrary{positioning,arrows.meta,calc,fit,patterns}
\usepackage[active,tightpage]{preview}
\PreviewEnvironment{tikzpicture}
\setlength\PreviewBorder{6pt}
% palette
\definecolor{cgen}{RGB}{60,60,60}
\definecolor{csim}{RGB}{0,114,178}
\definecolor{ccalo}{RGB}{213,94,0}
\definecolor{cpf}{RGB}{0,158,115}
\definecolor{creco}{RGB}{204,121,167}
\definecolor{cpu}{RGB}{150,150,150}
\definecolor{ctrack}{RGB}{86,180,233}
\definecolor{cnote}{RGB}{90,90,90}
\tikzset{
  ladderbox/.style={draw, rounded corners=3pt, thick, align=center, minimum width=3.0cm, minimum height=1.15cm, font=\small, fill=white},
  effect/.style={font=\footnotesize\itshape, text=cnote, align=center},
  arrow/.style={-{Stealth[length=2.5mm]}, thick},
  part/.style={circle, draw, thick, inner sep=1.2pt, minimum size=5mm, font=\scriptsize, fill=white},
  gen/.style={part, draw=cgen, fill=cgen!15},
  simp/.style={part, draw=csim, fill=csim!15},
  lost/.style={part, draw=cpu, dashed, text=cpu},
  jet/.style={draw, thick, rounded corners=6pt, dashed, inner sep=5pt},
  tbl/.style={draw, rounded corners=2pt, fill=white, font=\scriptsize\ttfamily, align=left, inner sep=4pt},
  note/.style={font=\scriptsize, text=cnote, align=left},
}

\begin{document}
\begin{tikzpicture}[node distance=0.5cm and 0.8cm,
  lvl/.style={font=\small\bfseries, text=cnote, anchor=east, align=right},
  obj/.style={draw, thick, rounded corners=3pt, fill=white, font=\scriptsize, align=center, inner sep=3pt, minimum height=0.6cm},
  cell/.style={draw, thick, minimum width=0.55cm, minimum height=0.55cm, inner sep=0pt, font=\tiny}]
  % ---------------------------------------------------------------- level 1: the jet
  \node[lvl] (l1) at (-0.3,0) {jet};
  \node[obj, draw=creco, minimum width=9.6cm, anchor=west] (jet) at (0,0) {reconstructed jet (PF, CHS or Puppi): anti-$k_T$ 0.4 of the candidates, each with its clustering weight $w_c$ (Puppi weight, 1 otherwise)};
  % ---------------------------------------------------------------- level 2: candidates
  \node[lvl] (l2) at (-0.3,-1.5) {candidates};
  \node[part, draw=creco, fill=creco!15] (c1) at (1.0,-1.5) {$h^\pm$};
  \node[part, draw=creco, fill=creco!15] (c2) at (3.6,-1.5) {$\gamma$};
  \node[part, draw=creco, fill=creco!15] (c3) at (6.4,-1.5) {$h^0$};
  \node[part, draw=cpu, fill=cpu!15] (c4) at (9.0,-1.5) {$h^\pm$};
  \foreach \c in {c1,c2,c3,c4} \draw[arrow, creco!70] (jet.south -| \c) -- (\c);
  % ---------------------------------------------------------------- level 3: block elements
  \node[lvl] (l3) at (-0.3,-3.2) {PF block\\ elements};
  \node[obj, draw=ctrack] (t1) at (1.0,-3.2) {track};
  \node[obj, draw=ccalo, dashed] (k1) at (2.3,-3.2) {HCAL\\ cluster};
  \node[obj, draw=ccalo] (k2) at (3.6,-3.2) {ECAL\\ cluster};
  \node[obj, draw=ccalo] (k3) at (5.4,-3.2) {ECAL\\ cluster};
  \node[obj, draw=ccalo] (k4) at (7.0,-3.2) {HCAL\\ cluster};
  \node[obj, draw=ctrack] (t4) at (9.0,-3.2) {track};
  \draw[arrow, ctrack] (c1) -- (t1);
  \draw[ccalo, dashed] (c1) -- (k1);
  \draw[arrow, ccalo] (c2) -- (k2);
  \draw[arrow, ccalo] (c3) -- (k3);
  \draw[arrow, ccalo] (c3) -- (k4);
  \draw[arrow, ctrack] (c4) -- (t4);
  \node[note, text width=2.3cm, align=center] at (2.3,-4.15) {linked, not used: the momentum is the track's};
  % ---------------------------------------------------------------- level 4: rechits
  \node[lvl] (l4) at (-0.3,-5.3) {rechits};
  \foreach \x/\n in {2.9/r1, 3.6/r2, 4.3/r3} \node[cell, draw=ccalo, fill=ccalo!12] (\n) at (\x,-5.3) {};
  \foreach \x/\n in {5.0/r4, 5.7/r5} \node[cell, draw=ccalo, fill=ccalo!12] (\n) at (\x,-5.3) {};
  \foreach \x/\n in {6.6/r6, 7.3/r7, 8.0/r8} \node[cell, draw=ccalo, fill=ccalo!12] (\n) at (\x,-5.3) {};
  \foreach \r in {r1,r2,r3} \draw[arrow, ccalo!70] (k2) -- (\r);
  \foreach \r in {r3,r4,r5} \draw[arrow, ccalo!70] (k3) -- (\r);
  \foreach \r in {r6,r7,r8} \draw[arrow, ccalo!70] (k4) -- (\r);
  \node[note, anchor=west, text width=3.3cm] at (8.6,-5.3) {cluster $=\sum$ rechits $\times$ fractions; a rechit shared by two clusters is split by its fractions};
  \node[font=\tiny, text=ccalo] at (4.3,-5.75) {shared};
  % ---------------------------------------------------------------- level 5: the cell and its simulated deposits
  \node[lvl] (l5) at (-0.3,-7.6) {the cell};
  \begin{scope}[shift={(2.6,-8.9)}]
    \draw[->, thick] (-0.2,0) -- (0,0) -- (0,2.6) node[above, font=\scriptsize] {energy};
    % rechit bar
    \fill[ccalo!25] (0.4,0) rectangle (1.2,2.1);
    \draw[thick, ccalo] (0.4,2.1) -- (1.2,2.1);
    \draw (0.4,0) rectangle (1.2,2.1);
    \node[font=\tiny, align=center] at (0.8,-0.3) {rechit\\ $E_{\mathrm{rec}}$};
    % sim deposits bar, per particle
    \fill[csim!60] (1.9,0) rectangle (2.7,0.9);
    \fill[csim!30] (1.9,0.9) rectangle (2.7,1.4);
    \fill[cpu!70] (1.9,1.4) rectangle (2.7,1.8);
    \draw (1.9,0) rectangle (2.7,1.8);
    \node[font=\tiny, align=center] at (2.3,-0.3) {sim deposits\\ $\times$ scale $s$};
    \node[font=\tiny, text=csim] at (3.2,0.45) {$p_1$};
    \node[font=\tiny, text=csim] at (3.2,1.15) {$p_2$};
    \node[font=\tiny, text=cpu] at (3.2,1.6) {pile-up};
    % attributed rechit
    \fill[csim!60] (4.2,0) rectangle (5.0,1.05);
    \fill[csim!30] (4.2,1.05) rectangle (5.0,1.63);
    \fill[cpu!70] (4.2,1.63) rectangle (5.0,1.8);
    \fill[pattern=north east lines, pattern color=cnote] (4.2,1.8) rectangle (5.0,2.1);
    \draw (4.2,0) rectangle (5.0,2.1);
    \node[font=\tiny, align=center] at (4.6,-0.3) {attributed};
    \node[font=\tiny, text=cnote, anchor=west] at (5.1,2.05) {residual};
    \node[font=\tiny, text=csim, anchor=west] at (5.1,0.5) {signal};
    \node[font=\tiny, text=cpu, anchor=west] at (5.1,1.55) {pile-up};
  \end{scope}
  \draw[arrow, ccalo!70] (r4) -- (3.4,-6.8);
  \node[tbl, anchor=north west, text width=5.2cm, font=\scriptsize] (formula) at (10.6,-6.6) {%
    per cell, per particle $p$:\\[2pt]
    $E_p^{\text{att}} = \min(E_{\mathrm{rec}},\, s\,E_{\mathrm{sim}}) \cdot \dfrac{E_p^{\mathrm{sim}}}{E_{\mathrm{sim}}}$\\[6pt]
    residual $= E_{\mathrm{rec}} - \min(E_{\mathrm{rec}},\, s\,E_{\mathrm{sim}})$\\[2pt]
    (noise, out-of-time pile-up, calibration)\\[4pt]
    $E_{\mathrm{sim}}$: signal + in-time pile-up deposits of the cell\\
    $s$: rechit scale of the sim deposits per subdetector};
  \node[note, below=0.15cm of formula, text width=5.2cm] {the sim deposits: the Geant4 hits of the cell per calorimeter-entrance particle (signal) and per pile-up particle; the rechit: the pulse fit of the digitised samples};
  % ---------------------------------------------------------------- track side
  \node[tbl, draw=ctrack, anchor=north, text width=2.4cm, font=\scriptsize, align=center] (tp1) at (1.0,-6.6) {track\\ $\downarrow$ $\chi^2$ association\\ TrackingParticle\\ $\downarrow$\\ SimTrack};
  \draw[arrow, ctrack] (t1) -- (tp1);
  \node[note, anchor=north east, text width=2.3cm, align=right] at (-0.3,-8.4) {the whole candidate to that particle: signal or pile-up by its event id, unattributed if no match};
  % ---------------------------------------------------------------- back up
  \node[obj, draw=csim, fill=csim!6, minimum width=16.2cm, anchor=west, text width=15.6cm] (up) at (0,-13.4) {\textbf{back up the chain}: cell shares $\times$ rechit fraction $\to$ cluster $\to$ candidate (\texttt{signalFraction}, \texttt{puFraction}) $\to$ jet: $\;\vec p_T^{\,\text{sig}} = \sum_c w_c f_c^{\text{sig}} \vec p_{T,c}$ (signal part), the rest pile-up and residual.\\ The particles' Geant4 ancestry gives the GenJet the jet is matched to, with the shares of its constituents as the shared energy.};
  \draw[arrow, csim] (tp1.south) -- (up.north -| tp1.south);
  \node[tbl, draw=cpu, anchor=north, text width=2.0cm, font=\scriptsize, align=center, text=cpu] (tp4) at (9.0,-6.6) {track\\ $\downarrow$\\ TrackingParticle of a pile-up interaction $\Rightarrow$ pile-up;\\ none $\Rightarrow$ unattributed};
  \draw[arrow, ctrack] (t4) -- (tp4);
  \draw[arrow, csim] (tp4.south) ++(0,-0.3) -- (up.north -| tp4.south);
  \draw[arrow, csim] (5.2,-11.6) -- (up.north -| 5.2,0);
\end{tikzpicture}
\end{document}
